
{
\footnotesize
\setcounter{utterance}{1}
\setlength{\tabcolsep}{0pt}
\begin{supertabular}{c@{$\;$}|p{.15\linewidth}@{}p{.15\linewidth}p{.15\linewidth}p{.15\linewidth}p{.15\linewidth}p{.15\linewidth}}   
  \# & $\;$\textbf{Player A} & \multicolumn{4}{c}{\textbf{Game Master}} & $\;\:$\textbf{Player B}\\
  \hline

  \theutterance \stepcounter{utterance} & &
  \multicolumn{4}{p{0.7\linewidth}}{
    \cellcolor[rgb]{0.9,0.9,0.9}{
      \makecell[{{p{\linewidth}}}]{
        \tt {\tiny [A$\langle$GM$\rangle$B]}
        <GAME DESCRIPTION> ...  \\
        \tt Available items: 1 patata, 1 guanto, 2 righelli, 1 fischietto, 2 schede.
      }
    }
  } & \\

  \theutterance \stepcounter{utterance} & &
  \multicolumn{4}{p{0.7\linewidth}}{
    \cellcolor[rgb]{0.9,0.9,0.9}{
      \makecell[{{p{\linewidth}}}]{
        \tt {\tiny [A$\langle$GM]}
        Player A value function: patata: 0, guanto: 1, righello: 3, fischietto: 1, scheda: 1.
      }
    }
  } & \\
 
  \theutterance \stepcounter{utterance} & &
  \multicolumn{4}{p{0.7\linewidth}}{
    \cellcolor[rgb]{0.9,0.9,0.9}{
      \makecell[{{p{\linewidth}}}]{
        \tt {\tiny [GM$\rangle$B]}
        Player B value function: patata: 7, guanto: 0, righello: 0, fischietto: 3, scheda: 0.
      }
    }
  } & \\

  \theutterance \stepcounter{utterance} &
  \multicolumn{4}{p{0.6\linewidth}}{
    \cellcolor[rgb]{0.8,1,0.9}{
      \makecell[{{p{\linewidth}}}]{
        \tt {\tiny [A$\rangle$B]}
        I miei valori: patata 0, guanto 1, righello 3, fischietto 1, scheda 1. Quali sono i tuoi valori? Collaboriamo per massimizzare il punteggio totale.
      }
    }
  } & & \\

  \theutterance \stepcounter{utterance} & & &
  \multicolumn{4}{p{0.6\linewidth}}{
    \cellcolor[rgb]{1,0.85,0.72}{
      \makecell[{{p{\linewidth}}}]{
        \tt {\tiny [GM$\langle$B]}
        {[}Proposta: 1 guanto, 2 righelli, 2 schede{]}
      }
    }
  } \\

  \theutterance \stepcounter{utterance} & &
  \multicolumn{4}{p{0.6\linewidth}}{
    \cellcolor[rgb]{0.9,0.9,0.9}{
      \makecell[{{p{\linewidth}}}]{
        \tt {\tiny [A$\langle$GM]}
        L'altro giocatore ha presentato una proposta segreta. Ora tocca a Lei presentare una sua proposta.
      }
    }
  } & \\

  \theutterance \stepcounter{utterance} &
  \multicolumn{4}{p{0.6\linewidth}}{
    \cellcolor[rgb]{0.8,1,0.9}{
      \makecell[{{p{\linewidth}}}]{
        \tt {\tiny [A$\rangle$GM]}
        {[}Proposta: 1 righello, 1 scheda{]}
      }
    }
  } & & \\
\end{supertabular}
}
